% Standalone mathematical module; requires amsmath, amssymb, amsthm, url.
% The surrounding document should supply theorem, lemma, and corollary environments.
% Written derivation, not a proof-assistant verification.
\section{Removing literal zero from the source of transport}
\label{zs:section}

The following is a restriction of a mathematical inference rule, rather
than a change in the spelling of zero. Throughout this section,
``nonzero source'' means \emph{not syntactically the single term $0$}.
Every variable and every term headed by $S$ or $I$ is a permitted source,
including terms whose numerical value is zero.

Let $\mathcal C$ be the native scalar calculus with terms
\[
 t::=x\mid0\mid St\mid I(n,a;u.h),
\]
the two computation schemas
\[
 I(0,a;u.h)=a,\qquad
 I(Sn,a;u.h)=h[I(n,a;u.h)/u],
\]
formulas generated by equality and the $Q,D$ binders, and the Hyp, QE,
DE, case-DI and equation-progress-DI rules stated in the main inventory.
The same argument applies to the original presentation with $D$ object
formulas and $Q$ resources in contexts. Let $\mathcal C^{-}$ replace
unrestricted transport by
\[
 \frac{\Gamma\vdash s=t\qquad\Gamma\vdash F(s)}
      {\Gamma\vdash F(t)}
 \quad\mathsf{Tr}^{-},\qquad s\not\equiv0.
 \tag{Tr$^-$}\label{zs:restricted-rule}
\]
Here $F$ is an arbitrary native numerical template, and $s\not\equiv0$
is a condition on finite term syntax. All other rules are unchanged.

\begin{theorem}[Literal-zero source elimination]
\label{zs:main}
For every native context $\Gamma$ and formula $A$,
\[
 \Gamma\vdash_{\mathcal C^{-}}A
 \quad\Longleftrightarrow\quad
 \Gamma\vdash_{\mathcal C}A.
\]
In particular, the ordinary whole-domain realization of
$\mathcal C^{-}$ is bi-interpretable and equiconsistent with full
first-order iPRA through the same interpretations as the native scalar
core. No change of numerical domain or interpreted equality is needed.
\end{theorem}

We prove the theorem without assuming arbitrary numerical substitution
admissibility for $\mathcal C^{-}$: such substitution can turn a variable
source into the forbidden literal zero. Every occurrence of a helper
equation at an arbitrary term is either obtained directly from its
displayed derived rule at that term, or from a derivation whose transport
sources retain an $S$ or $I$ head under the indicated substitution.

\subsection{Structural operations, cases, and equation induction}

Weakening, ordinary-assumption cut, and capture-free renaming are
admissible by induction on derivations. These operations do not turn a
nonzero source into literal zero. As usual, fresh parameters are renamed
before a cut. We also use the following restricted substitution fact:
a capture-free substitution is admissible for a particular finite
derivation if it leaves every primitive transport source in that
derivation syntactically different from zero. This fact follows by the
same induction; it is not unrestricted substitution admissibility.

\begin{lemma}[Cases and equation induction need no transport]
\label{zs:cases-induction}
In $\mathcal C^{-}$ the following rules are derivable:
\[
 \frac{\Gamma\vdash A(0)\qquad\Gamma\vdash A(Sz)}
      {\Gamma\vdash A(t)}\quad\mathsf{NC},
 \qquad
 \frac{\Gamma\vdash E(0)\qquad\Gamma,E(z)\vdash E(Sz)}
      {\Gamma\vdash E(t)}\quad\mathsf{EI}.
\]
In NC, $A$ is any native formula; in EI, $E$ is one literal equation.
The parameter $z$ is fresh for $\Gamma$, $t$, and the fixed parameters
of the matrix. Other parameters stay fixed throughout EI.
\end{lemma}
\begin{proof}
Choose an unused witness coordinate $x$. In NC, weaken the successor
proof by $A(0)$ and apply case-DI, with witness zero, to obtain
$D_{x,y}A(y)$. DE opens $Q_yA(y)$, because $x$ is unused. QE with the
given base proof and target $t$ yields $A(t)$, and DE closes.
For EI use progress-DI instead, with the given equation step, and use
the same DE/QE construction. No transport rule occurs.
\end{proof}

These rules will be applied to unconditional helper equations before
they are weakened into the context containing $0=t$. We never use NC
with a case parameter occurring in its retained context.

\subsection{The equality operations available immediately}

\begin{lemma}[Reflexivity and safe equality operations]
\label{zs:safe-equality}
Reflexivity $r=r$ is derivable for every term $r$. From $p=q$ with
$p\not\equiv0$, one can derive $q=p$ and $H(p)=H(q)$ for any numerical
template $H$. Every direct computation can be used in either direction
as an equation; using the reversed equation as a transport major is
permitted only when its new source is syntactically nonzero.
\end{lemma}
\begin{proof}
Put $b=I(0,r;u.0)$. The computation $b=r$, used as both major and
minor of transport with template $v=r$, gives $r=r$. Its source is
the nonempty iterator term $b$.
For symmetry of $p=q$, transport the reflexivity $p=p$ through the
template $v=p$. For numerical congruence, transport the reflexivity
$H(p)=H(p)$ through $H(p)=H(v)$. Both use the permitted source $p$.
Every computation equation has an $I$-headed left side, so its
symmetry follows by the same construction.
\end{proof}

An occurrence-selective template may keep some occurrences of a term
fixed while abstracting other occurrences. This is ordinary
capture-free template substitution, not an extra inference rule.

\begin{lemma}[Joining two nonzero terms known to be zero]
\label{zs:zero-join}
Suppose $p,q\not\equiv0$ and $\Gamma$ proves $p=0$ and $q=0$.
Then $\Gamma\vdash p=q$ using only $\mathsf{Tr}^{-}$ and computation.
\end{lemma}
\begin{proof}
Let
\[
 z_0=I(0,0;u.0),\qquad b_p=I(0,p;u.0).
\]
Computation gives $b_p=p$. Start with $b_p=b_p$ and use the major
$p=0$ in the template $I(0,v;u.0)=b_p$. This gives $z_0=b_p$.
Symmetry at the nonzero source $z_0$ gives $b_p=z_0$.
Now transport through $v=z_0$ using the computation $b_p=p$, whose
source is $b_p$, to obtain $p=z_0$.
The identical construction gives $q=z_0$. Symmetry at source $q$
gives $z_0=q$, and transport at source $z_0$ in $p=z_0$ gives $p=q$.
The sources used are $p,q,z_0,b_p,b_q$; all are syntactically nonzero.
No reverse transport from zero has been used.
\end{proof}

The zero-join lemma is deliberately placed before the next construction.
The positive case of its singleton test needs precisely this lemma;
ordinary transitivity through zero is not yet available.

\subsection{A singleton-one test without predecessor}

Define the finite scalar terms
\[
\begin{aligned}
 C(v,a,b)&:=I(v,a;u.b),& Z(v)&:=C(v,1,0),\\
 K(v)&:=C(v,0,1),& H(v)&:=C(v,1,2),\\
 N(n)&:=I(n,0;u.Z(u)),& R(n)&:=I(n,0;u.H(u)),\\
 L(n)&:=N(R(n)).
\end{aligned}
\tag{Tests}\label{zs:tests}
\]
The update variable in $C$ is unused. Here $1=S0$ and $2=S1$.
The standard-model description of $R$ is $0,1,2,2,\ldots$; the proof
below establishes the needed equations internally and does not assume
this description.

\begin{lemma}[Uniform test equations]
\label{zs:singleton}
The following equations have derivations in $\mathcal C^{-}$,
uniformly in the displayed argument:
\[
 H(v)=S K(v),\qquad R(1)=1,\qquad R(SSv)=2,
\]
\[
 N(1)=1,\quad N(2)=0,\quad
 L(1)=1,\quad L(SSv)=0,\quad
 \boxed{L(Sv)=Z(v)}.
 \tag{Singleton}\label{zs:singleton-equations}
\]
All transport sources in these helper derivations have an $I$ or $S$
head, except for the explicit nonzero endpoints in the zero-join
lemma; those endpoints also have an $I$ head in this application.
\end{lemma}
\begin{proof}
For $H(v)=S K(v)$ use NC. At zero both sides reduce to $1$; at a
successor both reduce to $2$. To join two equations $p=j$, $q=j$
with $j=1$ or $j=2$, reverse $q=j$ at its nonzero source $q$, then
transport at the nonzero numeral $j$. Each reduction of these
particular terms is a forward replacement of an iterator computation.
This proves the helper at an arbitrary target $v$ by NC, without
substituting zero into a variable-source transport.

The equations $N(1)=1$ and $N(2)=0$ are finite computations:
$N(0)=0$, $N(1)=Z(N(0))=1$, and $N(2)=Z(N(1))=0$.
Every replacement source in these calculations is an iterator term.
Likewise $R(1)=H(R(0))=H(0)=1$.

Computation and the first helper give, for any $v$,
\[
 R(Sv)=H(R(v))=S K(R(v)).
\]
Use this equation at the nonzero source $R(Sv)$ in
$R(SSv)=H(R(Sv))$. The resulting right side is
$H(S K(R(v)))=2$ by the successor computation of $C$.
Thus $R(SSv)=2$. These substitutions are all at $I$-headed sources.
Substitute $R(1)=1$ and $R(SSv)=2$ inside $N$ and use the two
computed values of $N$ to obtain $L(1)=1$ and $L(SSv)=0$.

Finally apply NC to the matrix $L(Sv)=Z(v)$.
At zero, $L(1)=1$ and $Z(0)=1$ are joined through the nonzero numeral
$1$. At a successor, $L(SSv)=0$ and $Z(Sv)=0$ are joined by
Lemma~\ref{zs:zero-join}; both left sides are nonempty iterator terms.
This proves the boxed equation at every target. Its use is therefore
an instance of a displayed derived rule, not an appeal to unrestricted
substitution admissibility.
\end{proof}

\begin{lemma}[An unconditional zero-filter identity]
\label{zs:filter}
For every term $t$,
\[
                 t=C(Z(t),t,0)
                 \tag{Filter}\label{zs:filter-equation}
\]
is derivable in the empty context of $\mathcal C^{-}$.
\end{lemma}
\begin{proof}
Apply NC to $v=C(Z(v),v,0)$, with the target $t$.
At zero, forward computations give
$C(Z(0),0,0)=C(1,0,0)=0$; reverse this equation at its nonzero
left side. At $Sv$, forward computations give
$C(Z(Sv),Sv,0)=C(0,Sv,0)=Sv$; again reverse at the nonzero
iterator left side. Thus the base and successor proofs have no
retained assumption depending on the case variable. All numerical
arguments of the final NC target are permitted. In particular, this
identity is proved before any assumption $0=t$ is used.
\end{proof}

\subsection{Turning a zero-source equality around}

\begin{lemma}[Zero-source symmetry]
\label{zs:zero-symmetry}
For every term $t$,
\[
                     0=t\ \vdash_{\mathcal C^{-}}\ t=0.
                     \tag{ZS}\label{zs:zero-symmetry-equation}
\]
\end{lemma}
\begin{proof}
If $t\equiv0$, use reflexivity. Otherwise $t$ is a permitted source.
Work in the fixed context $\Gamma=\{0=t\}$ and put
\[
                  A(n):=I(n,t;u.Su).
\]
We first prove $n=A(n)$ by EI, keeping $t$ fixed.

For the base, computation gives $A(0)=t$. Its symmetry gives
$t=A(0)$. Since $t\not\equiv0$, transport at source $t$ in the
retained assumption $0=t$ gives $0=A(0)$.

For the step, retain $z=A(z)$ with $z$ fresh. This source is a
variable, hence is permitted. Numerical successor congruence from
Lemma~\ref{zs:safe-equality} gives $Sz=S A(z)$.
Computation gives $A(Sz)=S A(z)$, and its symmetry gives
$S A(z)=A(Sz)$. Its source $S A(z)$ is also permitted, so transport
in $Sz=S A(z)$ gives $Sz=A(Sz)$. This is exactly the EI step.
The induction hypothesis is used at the same fixed parameter $t$.

At target $1$, EI gives $1=A(1)$. Forward computation gives
$A(1)=St$: the successor equation produces $S A(0)$, and the
nonzero-source computation $A(0)=t$ is used inside $S$.
Transport at source $A(1)$ consequently gives
\[
                              1=St.
\]
Its source $1$ is permitted. Numerical congruence for $L$ gives
$L(1)=L(St)$. Use $L(1)=1$ at source $L(1)$, and
$L(St)=Z(t)$ at source $L(St)$, obtaining $1=Z(t)$.
Symmetry at source $1$ gives $Z(t)=1$.

Weaken the unconditional Filter identity into $\Gamma$.
Transport at the iterator source $Z(t)$ changes its right side to
$C(1,t,0)$. The latter computes to zero, and its source is again
an iterator term. The resulting equation is $t=0$, as required.
\end{proof}

This proof uses cases only for unconditional helper equations. The
sole induction with the retained assumption $0=t$ has the literal
equation $n=A(n)$ as invariant. It does not perform a forbidden
case distinction while silently changing the retained context.

\subsection{Recovering equality algebra and formula transport}

\begin{lemma}[Full numerical equality algebra]
\label{zs:full-equality}
The calculus $\mathcal C^{-}$ admits unrestricted symmetry,
numerical congruence, and transitivity. In particular it admits
root replacement
\[
 \frac{\Gamma\vdash s=t\qquad\Gamma\vdash s=r}
      {\Gamma\vdash t=r}.
\]
\end{lemma}
\begin{proof}
Symmetry uses Lemma~\ref{zs:safe-equality} when the source is not
zero and Lemma~\ref{zs:zero-symmetry} otherwise.

For numerical congruence from $s=t$, a nonzero $s$ is handled by
Lemma~\ref{zs:safe-equality}. If $s\equiv t\equiv0$, use
reflexivity. If $s\equiv0$ and $t\not\equiv0$, first derive $t=0$,
apply permitted numerical congruence at source $t$ to obtain
$H(t)=H(0)$, and use the now established full symmetry.

For transitivity from $r=s$ and $s=t$, a nonzero middle term $s$
is handled by one permitted transport. If $s\equiv0$ and either
endpoint is literally zero, a given premise is already the
conclusion. Otherwise $r,t\not\equiv0$; reverse $0=t$ to get
$t=0$ and apply the zero-join lemma to $r=0$ and $t=0$.
Root replacement follows by symmetry and transitivity.
These constructions are acyclic: zero-source symmetry was proved
before any unrestricted numerical congruence or transitivity.
\end{proof}

\begin{lemma}[Full transport]
\label{zs:full-transport}
For every native numerical template $F(v)$, the rule
\[
 \frac{\Gamma\vdash s=t\qquad\Gamma\vdash F(s)}
      {\Gamma\vdash F(t)}
\]
is admissible in $\mathcal C^{-}$, without a source restriction.
\end{lemma}
\begin{proof}
Use external induction on the finite number of formula binders in
$F$, uniformly in retained contexts and fixed numerical parameters.
For an equation $F(v)\equiv A(v)=B(v)$, use unrestricted numerical
congruence to obtain $A(s)=A(t)$ and $B(s)=B(t)$. Join these with
$A(s)=B(s)$ using the established equality algebra.

For $F(v)=D_{x,y}B(v,x,y)$, eliminate $F(s)$ at a fresh $a$,
retaining $Q_yB(s,a,y)$. Introduce the desired $F(t)$ by case-DI
with witness $a$. Its local assumption is $B(t,a,0)$.
Full symmetry gives $t=s$, and the external induction hypothesis
for the smaller matrix transports the assumption to $B(s,a,0)$.
QE gives $B(s,a,Sz)$, and the induction hypothesis in the forward
direction gives $B(t,a,Sz)$. Close case-DI and DE. All new
parameters can be chosen fresh for the source equality and the
retained context.

If $Q$ is first-class formula syntax, first observe that its case
introduction is derived without transport: introduce
$D_{u,y}B(y)$ with unused $u$ by case-DI and eliminate it with the
one-node continuation $Q_yB(y)$. For $F(v)=Q_yB(v,y)$, use this
derived introduction for the target. Convert its base assumption
backwards by the induction hypothesis, apply the old Q resource,
and convert the successor conclusion forwards. This is the same
smaller-matrix argument and uses no new numerical induction.
\end{proof}

\begin{proof}[Proof of Theorem~\ref{zs:main}]
Every restricted derivation is an unrestricted one. Conversely,
replace each unrestricted transport node by the derivation supplied
by Lemma~\ref{zs:full-transport}, recursively translating its
premises first. Other inference nodes are unchanged. The result has
the same context and conclusion. Thus the full native consequence
relation, not just closed theorems, is preserved. The existing
native/ordinary and ordinary/iPRA interpretations and both identity
composites therefore transfer unchanged.
\end{proof}

\subsection{Compatibility with restricted computation certificates}
\label{zs:certificate-computation}

Some compact checkers require the initializer of a computation packet
to have nonempty syntax, and encode a successor counter by a nonempty
predecessor. These are certificate restrictions, not restrictions
on the two mathematical computation schemas. They do not introduce
a circular dependence on zero-source transport in the preceding proof.

Use the inherited certified-zero term
\[
                         c_0=I(2,1;u.0).
\]
Its equation $c_0=0$ is a directly admitted successor computation:
the predecessor and initializer are both $1$. For an arbitrary
computation instance with literal-zero initializer or literal-zero
successor predecessor, form its whole equation as a template $J(v)$,
replacing just these offending occurrences by $v$. The instance
$J(c_0)$ is a directly admitted computation. Transport through this
equation template using $c_0=0$ yields $J(0)$, the desired original
instance. Both computation premises are admitted directly, and the
only transport source is the nonempty term $c_0$.

This is exactly the inherited three-node normalization in
\path{Research/certified_zero/tests/producer.py}, function
\texttt{expand\_computation}: one transport and two computation
premises. Its source is $c_0$, even when both defects are corrected
simultaneously. Passive zero occurrences in an update are left alone.
Thus all computation instances used in the mathematical proof can
be expanded before compiling to the source-restricted checker.

\subsection{Scope of the result}

The new statement is the admissibility of all omitted literal-zero
source transports. The native term/formula language, arithmetic
schemas, and target theory remain the inherited ones. This theorem
justifies replacing a two-way equality/Q major-pattern test by a
source-discriminated test in a suitable certificate protocol. It
does not assert a minimum proof-checker size, a shortest proof
expansion, or a new characterization of the primitive recursive
functions. The displayed proof is a written universal derivation;
executable certificates, when supplied, validate particular
instances rather than substitute for this argument.

